\documentclass[11pt]{article}
\usepackage[margin=1in]{geometry}
\usepackage{amsmath,amssymb,mathtools}
\usepackage{booktabs}
\usepackage{array}
\usepackage{url}
\usepackage{xspace}
\usepackage[colorlinks=true,linkcolor=blue,citecolor=blue,urlcolor=blue]{hyperref}
\hypersetup{pdftitle={Mixed fast matrix multiplication: fused levels of accurate square and rectangular schemes over a vendor GEMM},pdfauthor={Greg Sidebottom}}
\newcommand{\R}{\mathbb{R}}
\newcommand{\mm}[1]{\langle #1\rangle}
\newcommand{\ginf}{\gamma_{\infty,2}}
\newcommand{\amp}{\operatorname{amp}}
\newcommand{\vect}{\operatorname{vec}}
\newcommand{\lv}[1]{\texttt{#1}}
\setlength{\emergencystretch}{1em}
\renewcommand{\topfraction}{0.85}
\renewcommand{\bottomfraction}{0.5}
\renewcommand{\textfraction}{0.15}
\renewcommand{\floatpagefraction}{0.75}

\title{Mixed fast matrix multiplication:\\
fused levels of accurate square and rectangular schemes\\
over a vendor GEMM, chosen by a calibrated planner}
\author{Greg Sidebottom}
\date{Draft --- October 28, 2026}

\begin{document}
\maketitle

\begin{abstract}
A fast matrix multiplication in practice is a few levels of a Strassen-like
scheme over a vendor GEMM.  We treat the levels as a menu: every level may use
a different scheme, square or rectangular, and a plan is a sequence of levels
ending in a GEMM leaf.  On Apple's M4~Pro, with Accelerate's GEMM at the
leaves, the menu is Strassen--Winograd $\mm{2\times2\times2:7}$, the
$262$-operation accurate $\mm{4\times4\times4:48}$ algorithm of Dumas, Pernet
and Sedoglavic, and their accurate $\mm{3\times3\times6:40}$ algorithm in its
three orientations, which bring a square problem back to square after three
levels.  A level is three streaming passes whose cost is the number of blocks
they move, so a scheme's worth is its saving of products per block moved and
its amplification of the leaf error; a screen of 70 published schemes by these
two measures leaves exactly this menu, and shows that nothing larger than
$4\times4$ and no Kronecker composite can help.  A calibrated dynamic program
chooses the plan per shape, precision and error budget; it prices GEMM streams
and pipelined passes that use the whole machine, and prices rectangular
problems from the square calibration until they are calibrated themselves.
At $n=13824$ the planner's mixed plans run $1.73\times$ (f32) and $1.91\times$
(f64) faster than Accelerate's GEMM, at $n=6912$ $1.38\times$ and
$1.71\times$.  Measured against the exact product, the $3\times3\times6$-family
plan that runs in the time of the best plan from the $4\times4$ scheme alone
has 23--43\% less error at every size and precision tried, and less than the
GEMM's own; the one exception is the planner's fastest f64 plan at $13824$,
which spends a Strassen--Winograd level on 12\% more speed at the price of
51\% more error than the best $4\times4$ plan.  As a test of the opposite
design -- cheap passes instead of accurate ones -- a ternary
$\mm{3\times3\times6:42}$ scheme from Perminov's catalogue, whose three
programs need 139 additions against about 370 operations for the accurate
scheme's, is measured as a level in the same library: its passes cost exactly
the same (they are bound by the blocks they move), and its two extra products
and $2.8\times$ amplification make it slower and less accurate than the
accurate level at every size tried.  Apple's GPU, through the vendor's
MPSGraph GEMM, ties the f32 plan at $13824$ in time with $2.6\times$ its
error -- its product is bit for bit Accelerate's -- and is $1.4$--$1.6\times$
faster below that size; it has no f64.  Two implementations of the library,
in C and in Rust, agree bit for bit.
\end{abstract}

\section{Introduction}\label{sec:intro}

A scheme for multiplying a $g_1\times g_2$ matrix by a $g_2\times g_3$ matrix
in $r<g_1g_2g_3$ products, applied to blocks, saves products at every level of
recursion.  On a real machine a level is not free: its linear combinations of
blocks are passes over memory, and the products it hands down are GEMMs whose
rate depends on their size and shape.  Libraries that use fast schemes
therefore apply one to three levels on top of a tuned GEMM and stop
\cite{HSvdG16,DPSST26}.  Which scheme, how many levels, in what order, and at
what cost in accuracy, are the questions this report answers for one machine,
Apple's M4~Pro with Accelerate's GEMM, by measurement.

The approach is to mix.  A \emph{plan} is a sequence of \emph{levels} from the
top of the problem down to a GEMM leaf, and every level may come from a
different scheme: Strassen--Winograd's $\mm{2\times2\times2:7}$ (\lv{SW}), the
accurate $\mm{4\times4\times4:48}$ algorithm of Dumas, Pernet and Sedoglavic
\cite{DPS25,DPS26a} in its $262$-operation form \cite{PLinOpt,S26b} (\lv{D48}),
and their accurate $\mm{3\times3\times6:40}$ algorithm \cite{DPS26b,PLinOpt} in
its three orientations $3\times3\times6$, $3\times6\times3$ and $6\times3\times3$
(\lv{R336}, \lv{R363}, \lv{R633}).  The three orientations applied in turn
divide a square problem by $54$ and leave it square, where three $4\times4$
levels divide by $64$.  A plan such as \lv{SW/R363+3/R633/G6} is one
Strassen--Winograd level, then a $3\times6\times3$ level whose products are
pipelined on three pass threads, then a $6\times3\times3$ level whose $40$
products run as six single-threaded GEMM streams.  Mixing pays for three
reasons that the measurements separate:

\begin{enumerate}
\item \emph{Order.}  A level writes more data than it reads (the $4\times4$
  scheme three times as much, Strassen's $1.75$ times); the larger expansion
  belongs at the smaller level, so \lv{SW/D48/G} moves $9.75\,n^2$ words per
  side where \lv{D48/SW/G} moves $12.25\,n^2$ with the same $336$ products.
\item \emph{Shape.}  A rectangular level hands down rectangular products; the
  $3\times3\times6$ cycle keeps $0.741$ of the products per level against the
  $4\times4$ scheme's $0.75$, moves $11.89\,n^2$ words against $12$, stays under
  $58$ streams per pass against $64$, and after two levels its products are
  larger, so Accelerate runs them 13\% faster.
\item \emph{Accuracy.}  Each level multiplies the leaf error by a factor set by
  the scheme and the point chosen in its orbit; the $3\times3\times6$
  algorithm's factor is $0.74$--$0.79$ of the $4\times4$ algorithm's per level in
  f32, and both are below one, because their products are shorter dot
  products than the GEMM's.  Under an error budget an accurate level admits
  deeper, faster plans.
\end{enumerate}

Two facts about the costs organize everything.  First, a fused pass is bound
by memory traffic, not arithmetic: the $201$-operation program we found for the
$4\times4$ scheme \cite{S26}, DPS's original $315$-operation program and their
$284$-operation accurate program all run at the same speed as fused passes
(totals within 1\%), so the operation counts that scheme searches minimize
have little value here, while the error growth of the program matters
\cite{S26b}.  Arithmetic shows only when a program is wasteful: the direct
sums of the $3\times3\times6$ scheme's $R$ and $P$ matrices take $18$
operations per output and their passes ran 30\% slower than the traffic
predicts, which PLinOpt's optimizer removed.  Second, the leaf GEMM dominates
the time (85--95\% of a plan), so a level pays only when its products are
large enough for the GEMM to run near its peak, and a scheme that saves
products beats one that saves additions.  The ternary $\mm{3\times3\times6:42}$
scheme of Section~\ref{sec:ternary} is the controlled experiment: a third of
the pass arithmetic, two products more, and it loses.

\paragraph{Contributions.}
(i)~A cost model of fused levels in blocks moved and GEMMs priced inside the
level that sits on them, calibrated per machine and precision, with
predictions within 0.3--3\% at $n\ge 8192$ (Section~\ref{sec:levels}).
(ii)~A screen of 70 published schemes by saving per block moved and by leaf
error amplification per $4\times4$-equivalent saving, which leaves the
accurate $\mm{3\times3\times6:40}$ cycle as the one addition to the
$4\times4$ menu and rules out every larger shape and every Kronecker composite
(Section~\ref{sec:menu}).  (iii)~A per-level error model, validated within
25\%, and the orbit floors of the amplification for the two accurate classes
(Section~\ref{sec:accuracy}).  (iv)~A planner over $(m,k,n)$ shapes with GEMM
streams, pipelines, padding to multiples of $2^a3^b$, an error budget, and
shape estimates that price uncalibrated rectangular problems within 2\%
(Section~\ref{sec:planner}).  (v)~Measurements on the M4~Pro:
$1.73\times$ and $1.91\times$ over Accelerate at $n=13824$ in f32 and f64,
with less error than the GEMM in both precisions for every plan without a
Strassen level and 23--43\% less than the best $4\times4$-only plan of the
same speed; rectangular problems at $1.04$--$1.52\times$; and the GPU as
the alternative, tied in time at $13824$ f32 with $2.6\times$ the error
(Section~\ref{sec:results}).  (vi)~The ternary trial
(Section~\ref{sec:ternary}).  (vii)~Two implementations, C and Rust, that
agree in plans, predicted times and results bit for bit
(Section~\ref{sec:impl}).  Code, programs, calibrations and raw results are
in the repository \cite{logic}.

\paragraph{The result in two tables.}
Plans are written from the top of the problem down to the leaves, the levels
separated by slashes.  \lv{D48} is a level of the accurate
$\mm{4\times4\times4:48}$ scheme: it cuts $A$, $B$ and $C$ into $4\times4$
grids of blocks and forms the product with $48$ block products instead of
$64$.  \lv{R336}, \lv{R363} and \lv{R633} are levels of the accurate
$\mm{3\times3\times6:40}$ scheme in its three orientations, named by the grids
they cut $A$ and $B$ into: \lv{R363} cuts $A$ $3\times6$ and $B$ $6\times3$
and uses $40$ block products instead of $54$.  \lv{SW} is a Strassen--Winograd
level ($2\times2$ grids, $7$ products instead of $8$).  \lv{G} is the leaf,
Accelerate's GEMM on the block products that remain, and \lv{G6} runs those
products as six concurrent single-threaded GEMM streams.  \lv{+p} after a
level runs its passes -- the sums that form the operands and accumulate the
outputs -- on $p$ threads overlapped with the products below.  So
\lv{R363+4/D48/G6} at $n=13824$ is a $3\times6\times3$ level with four pass
threads, whose $40$ products of $4608\times2304$ by $2304\times4608$ each get
a $4\times4$ level, whose $40\times48 = 1920$ leaf products of $1152\times576$
by $576\times1152$ go to Accelerate in six streams; the classical count would
be $54\times64 = 3456$.

The error of a product is the largest entry of $C - C_{\mathrm{exact}}$
divided by the largest entry of $C_{\mathrm{exact}}$, where
$C_{\mathrm{exact}}$ is the f64 product of the same inputs in f32 and a
compensated double-double product in f64, on uniform random inputs in
$[-1,1)$; every method in a comparison sees the same inputs.  The difference
from Accelerate's own result is not an error measure, because Accelerate's
error is as large as a plan's (Section~\ref{sec:results}).

Table~\ref{tab:intro-matched} is the accuracy claim: in each precision and
size, the $3\times3\times6$-family plan that runs in the time of the best
plan built from the $4\times4$ scheme alone has 23--43\% less error, and is
4--6\% faster besides.  Table~\ref{tab:intro-gpu} is the comparison with the
other engine on the machine, Apple's GPU through the vendor's MPSGraph GEMM,
in f32, the only precision it has.

\begin{table}[htb]
\centering\small
\begin{tabular}{llrrlrr}
\toprule
$n$, prec. & mixed plan & time (s) & error & best $4\times4$-only & time (s) & error\\
\midrule
6912 f32 & \lv{R363+4/D48/G6} & 0.141 & $2.26\cdot10^{-6}$ & \lv{D48+4/D48/G6} & 0.149 & $2.94\cdot10^{-6}$\\
13824 f32 & \lv{R363+4/R363/G6} & 0.907 & $2.46\cdot10^{-6}$ & \lv{D48+4/D48/G6} & 0.949 & $4.35\cdot10^{-6}$\\
6912 f64 & \lv{R363+3/R633/G6} & 0.451 & $3.75\cdot10^{-15}$ & \lv{D48+4/D48/G6} & 0.478 & $5.60\cdot10^{-15}$\\
13824 f64 & \lv{R363+3/D48/G6} & 3.467 & $5.81\cdot10^{-15}$ & \lv{D48+4/D48/G6} & 3.703 & $8.28\cdot10^{-15}$\\
\bottomrule
\end{tabular}
\caption{Matched time: the $3\times3\times6$-family plan against the best
plan from the $4\times4$ scheme alone (M4~Pro, square problems, best of
three to five rounds, error against the exact product).  Accelerate alone
takes $0.195$, $1.565$, $0.773$ and $6.223$~s on these four with errors
$3.85\cdot10^{-6}$, $6.40\cdot10^{-6}$, $6.69\cdot10^{-15}$ and
$1.23\cdot10^{-14}$.  The first three mixed plans are the planner's
choices; at $13824$ f64 the planner's choice is the faster
\lv{SW/R363+3/R633/G6} ($3.251$~s), whose error $1.25\cdot10^{-14}$ is the
GEMM's own (Section~\ref{sec:results}).}
\label{tab:intro-matched}
\end{table}

\begin{table}[htb]
\centering\small
\begin{tabular}{lrrrrrr}
\toprule
 & \multicolumn{2}{c}{GPU (MPSGraph)} & \multicolumn{2}{c}{Accelerate} & \multicolumn{2}{c}{best mixed plan}\\
\cmidrule(lr){2-3}\cmidrule(lr){4-5}\cmidrule(lr){6-7}
$n$, f32 & time (s) & error & time (s) & error & time (s) & error\\
\midrule
3456 & 0.013 & $3.04\cdot10^{-6}$ & 0.025 & $3.04\cdot10^{-6}$ & 0.021 & $1.28\cdot10^{-6}$\\
6912 & 0.099 & $3.85\cdot10^{-6}$ & 0.195 & $3.85\cdot10^{-6}$ & 0.141 & $2.26\cdot10^{-6}$\\
13824 & 0.946 & $6.40\cdot10^{-6}$ & 1.565 & $6.40\cdot10^{-6}$ & 0.907 & $2.46\cdot10^{-6}$\\
\bottomrule
\end{tabular}
\caption{The GPU against the CPU in f32 (the GPU has no f64): the M4~Pro's
20-core GPU through MPSGraph's matrix multiplication on zero-copy shared
buffers, wall clock; Accelerate's GEMM; and the best mixed plan (\lv{R363/G6},
\lv{R363+4/D48/G6}, \lv{R363+4/R363/G6}).  The GPU's product is bit for bit
Accelerate's, so its error is the GEMM's (Section~\ref{sec:gpu}).}
\label{tab:intro-gpu}
\end{table}

\paragraph{Setting.}
Row-major $A\in\R^{m\times k}$, $B\in\R^{k\times n}$, $C=AB$, in f32 or f64.
The machine is an Apple M4~Pro (two performance clusters of five cores, one
efficiency cluster of four, 64~GB) running macOS~26; Accelerate's GEMM runs on
the SME matrix units, one per cluster, at about $3.3$~TFLOP/s in f32 and
$0.85$ in f64 on large square problems.  Times are the best of five
interleaved rounds (each round runs every method once; medians are within 2\%
unless stated).  Errors are measured as defined above (``err64'' in the
tables); the tables of Section~\ref{sec:results} also list ``relerr'', the
same max-norm difference taken against Accelerate's product, which includes
the GEMM's own error and is kept only for comparison with the earlier
campaigns.

\section{Fused levels and what they cost}\label{sec:levels}

\subsection{Schemes as linear programs}

A scheme $\mm{g_1\times g_2\times g_3:r}$ is three matrices $L\in\R^{r\times
g_1g_2}$, $R\in\R^{r\times g_2g_3}$, $P\in\R^{g_1g_3\times r}$ with
\begin{equation}\label{eq:scheme}
\vect(AB) \;=\; P\bigl((L\,\vect A)\circ(R\,\vect B)\bigr)
\qquad\text{for all } A\in\R^{g_1\times g_2},\ B\in\R^{g_2\times g_3},
\end{equation}
$\circ$ the entrywise product and $\vect$ row-major.  Applied to blocks, the
same identity holds with $A$'s $g_1g_2$ blocks of size $\frac{m}{g_1}\times
\frac{k}{g_2}$, $B$'s $g_2g_3$ blocks of $\frac{k}{g_2}\times\frac{n}{g_3}$,
and $C$'s $g_1g_3$ blocks of $\frac{m}{g_1}\times\frac{n}{g_3}$: the $r$
products are $(\frac{m}{g_1}\times\frac{k}{g_2})(\frac{k}{g_2}\times
\frac{n}{g_3})$ matrix products, computed by the plan below.  A \emph{level} is
the application of \eqref{eq:scheme} once; the matrices are realized as
straight-line programs (PLinOpt's \texttt{.slp} syntax \cite{PLinOpt}), one
per side, and the three programs are what a level executes:
\begin{itemize}
\item the $L$ pass reads the $g_1g_2$ blocks of $A$ and writes the $r$ left
  factors $S_i$;
\item the $R$ pass reads the $g_2g_3$ blocks of $B$ and writes the $r$ right
  factors $T_i$;
\item the products $M_i=S_iT_i$ run (GEMMs, or a level below);
\item the $P$ pass reads the $r$ products and writes the $g_1g_3$ blocks of $C$.
\end{itemize}
Each pass is \emph{fused}: one loop over the element positions of a block,
with the whole program evaluated on scalars in registers at each position.
A pass therefore reads every input block once and writes every output block
once, and the program's operation count costs only arithmetic, which the
memory system hides.

\subsection{The cost of a level}

Write $w$ for the word size and $B$ for the streaming bandwidth the passes
reach.  A level on an $(m\times k)(k\times n)$ problem moves
\begin{equation}\label{eq:traffic}
W(m,k,n) \;=\; (g_1g_2+r)\,\frac{mk}{g_1g_2} \;+\; (g_2g_3+r)\,\frac{kn}{g_2g_3}
\;+\; (g_1g_3+r)\,\frac{mn}{g_1g_3}
\end{equation}
words through its three passes, that is $n^2\bigl(3 + r(\tfrac{1}{g_1g_2} +
\tfrac{1}{g_2g_3} + \tfrac{1}{g_1g_3})\bigr)$ on a square problem, and costs
\begin{equation}\label{eq:level}
T(m,k,n) \;=\; r\,T_{\text{below}}\!\Bigl(\frac{m}{g_1},\frac{k}{g_2},
\frac{n}{g_3}\Bigr) \;+\; \frac{w\,W(m,k,n)}{B}.
\end{equation}
Two numbers of a scheme decide the second term: its \emph{traffic}
$W/n^2$ and its \emph{streams}, the number of blocks a pass reads and writes
at once, $g_1g_2+r$, $g_2g_3+r$ and $g_1g_3+r$.  Beyond about $64$ streams a
pass loses bandwidth: a $\mm{5\times5\times5:93}$ level's $118$-stream passes
ran at half the bandwidth of the $64$-stream passes of the $4\times4$ scheme.
Table~\ref{tab:levels} gives these numbers for the menu.  The first term is
decided by the \emph{saving} $r/(g_1g_2g_3)$ and by the GEMM's rate on the
products' shape and size: the f32 GEMM keeps $3.2$--$3.3$~TFLOP/s down to
$2048$ and drops to $3.0$ at $1024$ and $2.6$ at $512$, where the f64 GEMM
keeps its $0.81$--$0.85$~TFLOP/s down to $1024$, so f64 rewards depth and f32
does not.

\begin{table}[t]
\centering\footnotesize
\begin{tabular}{llrrrll}
\toprule
level & scheme & saving & streams $L/R/P$ & traffic ($n^2$) & amp$_{\text{eq}}$ & $\gamma_{\text{eq}}$\\
\midrule
\lv{SW} & Strassen--Winograd $\mm{2\times2\times2:7}$ & 0.875 & 11/11/11 & 8.25 & 10.66 & 19.8\\
\lv{D48} & DPS accurate $\mm{4\times4\times4:48}$, 262-op.\ point & 0.750 & 64/64/64 & 12.00 & 2.828 & 5.657\\
\lv{R336} & DPS accurate $\mm{3\times3\times6:40}$ & 0.741 & 49/58/58 & 11.89 & 2.280 & 6.162\\
\lv{R363} & the same, orientation $3\times6\times3$ & 0.741 & 58/58/49 & 11.89 & 2.453 & 5.125\\
\lv{R633} & the same, orientation $6\times3\times3$ & 0.741 & 58/49/58 & 11.89 & 2.280 & 6.162\\
\midrule
\lv{T336} & ternary $\mm{3\times3\times6:42}$ (Section~\ref{sec:ternary}) & 0.778 & 51/60/60 & 12.33 & 6.407 & 11.07\\
\lv{T363} & orientation $3\times6\times3$ & 0.778 & 60/60/51 & 12.33 & 6.166 & 9.737\\
\lv{T633} & orientation $6\times3\times3$ & 0.778 & 60/51/60 & 12.33 & 8.061 & 22.84\\
\bottomrule
\end{tabular}
\caption{The levels: saving of products per level, blocks moved by each pass,
pass traffic on a square problem in units of $n^2$ words, and the two
accuracy measures of Section~\ref{sec:menu}, normalized to the classical
algorithm of the shape and raised to the power $\ln0.75/\ln(\text{saving})$ so
that they are comparable per $4\times4$-equivalent saving.}
\label{tab:levels}
\end{table}

\subsection{What the implementation taught}

The passes reach $120$--$135$~GB/s against a peak of about $270$.  Two
details were each worth a large factor.  With $64$ streams the compiler gave
up vectorizing the fused loop (it could not rule out overlap among $64$
pointers); a \texttt{vectorize(assume\_safety)} pragma gave $4.3\times$ on the
passes.  And scratch blocks of power-of-two size all start at the same offset
modulo the cache's set span, so element $c$ of all $64$ streams fell into one
cache set; padding each block by $1040$ elements gave another $2.5\times$.
Operation counts did not matter: our $201$-operation program for the
$4\times4$ scheme \cite{S26}, DPS's $315$-operation original and their
$284$-operation accurate program run at the same speed as fused passes
(totals within 1\%; the $315$-operation passes take 2--8\% longer, about
0.3\% of the total).  The exception is a wasteful program: the direct sums of
the $3\times3\times6$ scheme's $R$ and $P$ matrices cost $18$ operations per
output, and their passes ran at $0.092$~s against $0.070$ for the $4\times4$
level's at $13824$ in f32 for $0.99\times$ the words.  PLinOpt's optimizer
\cite{PLinOpt} brought the nine programs of the three orientations from
$152$, $344{+}384$ and $366{+}384$ operations ($L$, $R$, $P$ of the $3\times3\times6$
orientation; additions plus scalings) to $59$, $117{+}22$ and $148{+}18$, which
cut the $R$-level passes by a quarter and the plans by 1--7\% with the error
unchanged.

\subsection{Using the whole machine}\label{sec:machine}

Accelerate drives each performance cluster's SME unit from one thread, so its
GEMM keeps two or three cores busy and leaves the rest idle.  The $r$ products
of a level are independent, and a leaf \lv{G}$k$ runs them as $k$
single-threaded GEMM streams taking products in turn: three streams beat
Accelerate's own threading by 10\% (f32) and 8\% (f64) at $2048$, more streams
win for the smallest products (f32 at $1024$, f64 at $512$) and lose badly
for large ones (six streams $+26$--$34$\% in f64 at $1024$).  The calibration
times each level's products with $0$, $3$, $4$ and $6$ streams and the
planner chooses per plan: 6--14\% off the best plans.

A streaming pass beside the SME streams slows them, but less than the pass
costs run alone: at f64 blocks of $512$ two pass threads slow six streams by
14\% while the pass runs at $49$~GB/s, about 35\% of its stand-alone cost.  So
a level whose products are themselves levels over GEMM streams can
\emph{pipeline} them (written \lv{X+p}, $p$ pass threads): while the streams
compute sub-problem $q$, the pass threads write sub-problem $q-1$'s result
and the factors of $q+1$, the scratch below double-buffered.  Pipelining
takes 8--22\% off the plans over six-stream leaves; one pass thread is too
few, and sub-problems under $1$~ms are not pipelined.  The calibration times
a pipeline of $2$ and of $10$ synthetic sub-problems per level, block size,
stream count ($3$, $6$) and pass threads ($2$, $3$, $4$): the slope is the
steady time per sub-problem, the rest the fill and drain.  NEON GEMMs on the
idle cores were tried and not adopted: at most 1--3\%, fragile, and the SME
units lose most of what the vector units add.

\section{The menu: which schemes can help}\label{sec:menu}

In a fused-pass library a plan already is a combination of sub-schemes, one
per level on the current $(m,k,n)$ sub-problem, so the question of mixing is
the question of the menu.  We scored every scheme on disk -- AlphaEvolve's 15
shapes \cite{Novikov25}, the 29 rectangular schemes of Kauers and Wood
\cite{KW} and the $5\times5\times6$ and $6\times6\times6$ schemes of
Moosbauer, Poole and Kemper \cite{MP25}, both as distributed with the
addition-reduction library of M\r{a}rtensson and Stankovski Wagner
\cite{MSW25}, PLinOpt's small and rectangular schemes including DPS's
accurate variants \cite{PLinOpt}, Smirnov's $\mm{3\times3\times6:40}$
\cite{Smirnov13}, and Kronecker composites of the accurate pieces -- by the
quantities that decide a level's cost and error in \eqref{eq:level}:
\begin{itemize}
\item saving $r/(g_1g_2g_3)$, streams $g_1g_2+r$, $g_2g_3+r$, $g_1g_3+r$, and
  traffic $3 + r\sum 1/(g_ig_j)$;
\item the amplification of the leaf error,
\begin{equation}\label{eq:amp}
\amp \;=\; \max_j\ \tfrac12\sqrt{\textstyle\sum_i P_{ji}^2\,\|L_i\|_2^2\,\|R_i\|_2^2},
\end{equation}
  the root-sum-square growth of independent rounding errors in the $r$
  products through output $j$, normalized by the classical algorithm's
  $\tfrac12\sqrt{g_2}$; and DPS's growth factor $\ginf=\max_j\sum_i
  |P_{ji}|\,\|L_i\|_2\|R_i\|_2$ \cite{DPS24,DPS26a}, normalized by $g_2$.
\end{itemize}
Schemes of different shapes save different fractions per level, so the
normalized factors are raised to the power $\ln0.75/\ln(\text{saving})$, the
growth per $4\times4$-equivalent saving, and are then directly comparable with
the $4\times4$ scheme's $2.828$ and $5.657$ (the columns amp$_{\text{eq}}$ and
$\gamma_{\text{eq}}$ of Table~\ref{tab:levels}).  The complete table is
\texttt{matmul/levels/levels.csv} \cite{logic}; Table~\ref{tab:menu} shows the
candidates.

\begin{table}[t]
\centering\footnotesize
\begin{tabular}{lrrrrl}
\toprule
scheme & saving & streams & traffic & amp$_{\text{eq}}$ & $\gamma_{\text{eq}}$\\
\midrule
$\mm{3\times3\times6:40}$ DPS accurate (and $6\times3\times3$) & \textbf{0.741} & 49/58/58 & \textbf{11.89} & \textbf{2.280} & 6.162\\
$\mm{3\times6\times3:40}$ DPS accurate & 0.741 & 58/58/49 & 11.89 & 2.453 & \textbf{5.125}\\
$\mm{4\times4\times4:48}$ accurate, 262-op.\ point & 0.750 & 64/64/64 & 12.00 & 2.828 & 5.657\\
$\mm{3\times4\times6:54}$ AlphaEvolve ($\pm1$ point) & 0.750 & 66/78/72 & 12.75 & 4.637 & 6.784\\
$\mm{3\times4\times7:63}$ rational & 0.750 & 75/91/84 & 13.50 & 5.734 & 11.13\\
$\mm{5\times5\times5:93}$ AlphaEvolve & 0.744 & 118/118/118 & 14.16 & 8.560 & 17.46\\
$\mm{5\times5\times6:110}$ Kemper & 0.733 & 135/140/140 & 14.73 & 10.38 & 22.9\\
$\mm{6\times6\times6:153}$ Moosbauer--Poole & 0.708 & 189/189/189 & 15.75 & 9.229 & 18.3\\
$\mm{2\times4\times5:32}$ AlphaEvolve & 0.800 & 40/52/42 & 11.80 & 4.962 & 10.95\\
$\mm{2\times2\times2:7}$ DPS accurate (in $\mathbb{Q}(\sqrt3)$) & 0.875 & 11/11/11 & 8.25 & 3.855 & 10.53\\
$\mm{2\times2\times2:7}$ Strassen / Winograd & 0.875 & 11/11/11 & 8.25 & 6.89 / 10.66 & 14.1 / 19.8\\
$\mm{3\times3\times3:23}$ best of four & 0.852 & 32/32/32 & 10.67 & 12.7 & 34.1\\
$\mm{3\times3\times6:40}$ Smirnov ($\pm1$ point, same class) & 0.741 & 49/58/58 & 11.89 & 9.754 & 26.1\\
\bottomrule
\end{tabular}
\caption{The menu screen (\texttt{matmul/levels/levels.py}): the schemes that
could matter, at their published points.  The accuracy columns are per
$4\times4$-equivalent saving.}
\label{tab:menu}
\end{table}

Three conclusions.  Everything larger than $4\times4$ saves more per level but
needs $110$--$225$ streams per pass, where the passes run at half bandwidth,
and has $8$--$100\times$ the classical amplification at its published points.
Kronecker products add nothing: a composite keeps the per-saving
amplification of its pieces ($\mm{3\times3\times6}\otimes\mm{3\times3\times6}$
keeps $2.280$, the accurate $2\times2$ squared keeps $3.855$) and multiplies
their stream counts, so it is never better than its two levels run one after
the other and the stream explosion makes it worse.  And the one scheme that
beats the $4\times4$ scheme on every count but $\gamma_{\text{eq}}$ is DPS's
accurate $\mm{3\times3\times6:40}$: applied to a square problem, the
$3\times3\times6$ orientation hands down $(n/3,n/3,n/6)$ products, the
$3\times6\times3$ orientation on those gives $(n/9,n/18,n/18)$, and the
$6\times3\times3$ orientation returns to the square $(n/54)^3$.  The cycle
keeps $0.741$ of the products per level, moves $11.89\,n^2$ words, stays
within $58$ streams, and amplifies the leaf error by $2.28$--$2.45$ per
$4\times4$-equivalent saving against $2.83$.  Perminov's catalogue of
$680$ shapes up to $16\times16\times16$ \cite{Perminov} confirms the menu: only
$\mm{3\times3\times6:40}$, $\mm{4\times4\times4:48}$ and $\mm{3\times4\times6:54}$
save a quarter of the products within $80$ streams per pass, and the new
shapes that save a quarter, $\mm{3\times6\times8:108}$ and $\mm{4\times7\times9:187}$,
need $126$--$215$ streams.

The paper cost model \eqref{eq:level} with the GEMM at the lone $8192^3$ f32
rate and the passes at the $4\times4$ level's bandwidth puts the cycle 1--3
points ahead of the $4\times4$ plans at every depth (\lv{R336/G} 17.8\% off the
GEMM against \lv{D48/G} 16.8\%; \lv{R336/R363/G} 13.9\% against 10.8\%), and
the assumption behind it holds: Accelerate on the rectangular leaves at equal
flops runs within 5\% of the square GEMM after one level and 13\% ahead after
two, because the leaves are larger ($n/54$ against $n/64$ after three levels).

\section{Accuracy}\label{sec:accuracy}

\subsection{What a level does to the error}

The rounding errors of the $r$ leaf products pass through $P$ into the
outputs; the roundings of the passes themselves add to them.  Taking the leaf
errors as independent and proportional to the sizes of the factors gives
\eqref{eq:amp}: a level multiplies the leaf error by $\amp$, against
$\tfrac12\sqrt{g_2}$ for the classical algorithm on the same blocks.  This
root-sum-square measure, not the worst-case sum $\ginf$, is what the
measurements follow: the $4\times4$ orbit's $\ginf$ minimum ($21.41$, in
$\mathbb{Q}(\sqrt2)$) has $\amp=3.14$ and measures no better than DPS's
Table~1 point, while the $262$-operation point, with $\amp=2\sqrt2$ the
orbit's minimum and the smallest $\ginf$ among dyadic points, measures
$0.75$--$0.77$ of it \cite{S26b}.

Over its whole orbit under $GL_3\times GL_3\times GL_6$ the
$\mm{3\times3\times6:40}$ class has amplification ratio at least $2.242$ in
every orientation (L-BFGS from random starts, \texttt{orbit\_rect.py}); DPS's
$3\times3\times6$ point is 5\% above that floor and 5\% above its $\ginf$
floor, their $3\times6\times3$ point 14\% above (and 41\% above that
orientation's $\ginf$ floor of $3.90$, far below the $4\times4$ class's
$5.66$).  The $\pm1$ Smirnov point of the same class (amplification ratio
$10.8$) reaches the same floors, so the class carries them, not the point.
Per $4\times4$-equivalent saving the floor is $2.17$ against the $4\times4$
class's $2.83$, which is its orbit's minimum and, as far as continuous search
from many starts can tell, a local minimum of the whole variety of rank-48
decompositions.  A better exact point for the $3\times6\times3$ orientation is
worth a dyadic search.

\subsection{The error model and its validation}

A plan's error is well described by the classical error scale times one
factor per fused level, growing like $\sqrt{k}$ with the inner dimension,
\begin{equation}\label{eq:errmodel}
\text{err}(\text{plan};\,m,k,n) \;=\; e_0\,\sqrt{k/N_E}\ \prod_{\text{levels }X} f_X,
\end{equation}
with the order of the levels immaterial (measured: \lv{SW} then \lv{D48} and
\lv{D48} then \lv{SW} within 2\%) and streams and pipelines leaving the
arithmetic unchanged.  The calibration measures six plans at $N_E=2304$
(divisible by $16$, $3$ and $6$) against the classical GEMM and fits $e_0$ and
the factors: in f32 \lv{SW} $1.09$, \lv{D48} $0.91$, \lv{R336} $0.74$, \lv{R363}
$0.79$, \lv{R633} $0.76$; in f64 $1.99$, $1.20$, $1.25$, $1.27$, $1.42$.  The
predictions hold within 25\% from $4096$ to $16384$ except for two Strassen
levels, over-predicted by up to 50\% (the safe side).  The factors are fit
to differences from Accelerate's product, which include the GEMM's own error,
as large as an accurate level's (Table~\ref{tab:results}); they separate the
Strassen level from the accurate ones but not the accurate levels from each
other -- in f64 they put \lv{R363} ($1.27$) behind \lv{D48} ($1.20$) where
the exact error puts it 36--41\% ahead -- and the budget \lv{auto:}$g$
inherits this.  The fit should move to the exact product (the compensated
product costs $0.1$~s at $N_E$); it has not yet.  The accurate levels'
factors are below one in f32: against the GEMM their products are $k/g_2$-long
dot products instead of $k$-long ones, and the GEMM's own accumulation error
is the larger term.  Measured against the exact product, one or two accurate
levels have less error than Accelerate's GEMM in both precisions
(Table~\ref{tab:results}).
The budget (\lv{max\_growth}, \lv{auto:}$g$; defined in
Section~\ref{sec:planner}) takes the fastest plan whose predicted growth
$\prod f_X$ is at most $g$; with these factors every plan the planner likes
qualifies down to $g=2$, where the first $4\times4$ program we used, our
$201$-operation point with factor $3.4$ per level, forced the classical blocked
GEMM at $g=2$ and one level at $g=8$.

\subsection{Programs and precision}

For a fixed scheme the program still matters, through the roundings of the
$P$ pass: a first-order rounding model with exact second moments of every
intermediate shows that DPS's $262$-operation programs collect $15$ roundings
of RMS size $0.80$ of the output per output where our $445$-operation
programs for the same point collect $36$ of size $1.39$, and among four
programs of equal amplification the model's worst-output error orders the
measured errors exactly at three and five levels \cite{S26b}.  The lever is
nearly exhausted for the $4\times4$ class: an 11\% smaller modelled pass
error was invisible at one to three levels and worth about 5\% at depth.  The
precision lever remains: running the passes of the f32 path in f64 (the
blocks stay f32, the arithmetic is f64, one rounding at the GEMM interface)
costs $\le2$\% below $16384$ and cuts the error by 36\% at five levels, nothing
at one to three, where the leaf GEMMs' error dominates.  It is a switch in the
library (\texttt{FMM\_F64PASS}), not the default.

\section{The planner}\label{sec:planner}

A plan is a sequence of levels from the top, ending in a leaf: \lv{G} (one
GEMM, Accelerate's threading), \lv{G}$k$ ($k$ GEMM streams), or the classical
blocked GEMM \lv{B2} on its own.  The planner is a dynamic program over the
first level:
\begin{equation}\label{eq:dp}
\text{cost}(s) \;=\; \min\Bigl(\text{leaf}(s),\ \min_X\bigl[\text{passes}_X(s)
+ r_X\,\text{cost}_X(s/g_X)\bigr]\Bigr),
\end{equation}
where $s=(m,k,n)$, $s/g_X$ is the shape of $X$'s products, and
$\text{cost}_X(t)$ is the cheaper of a GEMM leaf \emph{priced inside level
$X$} (with $0$, $3$, $4$ or $6$ streams) and a further level on $t$.  Pricing
the leaf inside the level matters: a GEMM shorter than about $1$~ms takes
13--27\% longer called back to back, as inside a plan, than called alone, and
the first calibration, which timed single calls, picked a two-level plan at
$8192$ f64 that ran 12\% slower than the right one.  A level whose
sub-problems are a level over GEMM streams may instead be pipelined
(Section~\ref{sec:machine}), priced from the calibrated steady time per
sub-problem plus fill and drain.

\paragraph{The error budget.}\label{par:budget}
Every call takes a number $g$: \lv{max\_growth} in the C and Rust
interfaces, \lv{auto:}$g$ in the benchmark.  For each candidate plan the
planner computes its predicted \emph{error growth}, the product
$\prod_X f_X$ of the error model \eqref{eq:errmodel}: the plan's predicted
error divided by the classical GEMM's predicted error on the same problem,
so that a growth of $1$ means as accurate as the GEMM, $0.5$ twice as
accurate, $2$ twice the error.  With $g>0$ the planner returns the fastest
plan whose growth is at most $g$: $g=1$ asks for a plan no less accurate than
the GEMM it replaces, $g<1$ for a more accurate one, and $g>1$ admits plans
up to $g$ times the GEMM's error, which is where the plans with a Strassen
level and the deep plans live.  $g=0$ is not ``no growth'' but ``no budget'':
the fastest plan whatever its error, and the setting of every timing in this
report unless a budget is named.  So $g=0$ can return a less accurate plan
than a positive $g$ would: at $13824$ f64 it returns \lv{SW/R363+3/R633/G6},
whose exact error is $1.02\times$ the GEMM's, where $g=1$ with a correctly
fit model would return \lv{R363+3/D48/G6} at $0.47\times$
(Table~\ref{tab:results}).  Two cautions.  The growth is the model's
estimate, good to about 25\% (Section~\ref{sec:accuracy}), so a budget is a
target rather than a guarantee, and near a boundary the choice is not
robust.  And the factors of the present calibration are fit against
Accelerate's result, which includes the GEMM's own error: in f64 they put
every level above $1$ (one $4\times4$ level at $1.2$ where the exact figure
is $0.63$), so until the model is refit to the exact product a budget of $1$
in f64 returns the GEMM and the budgets that admit the accurate levels are
$1.5$--$2$.  After the refit the sensible default is $g=1$, with $0$ for the
users who want the Strassen plans' speed.

\paragraph{Calibration.}
Everything the model needs is measured once per machine and precision and
cached (format~5, \texttt{calib5\_f32.txt}, \texttt{calib5\_f64.txt}, keyed by
the machine, the padding, the file format and the programs the kernels came
from): the lone GEMM and the blocked GEMM at each size; for each level and
product shape its three passes and the GEMM time per product at $0$, $3$, $4$
and $6$ streams, measured together by running the level; the pipelines; and
the six error-model plans at $N_E=2304$.  For $n$ up to $13824$ this is $408$
entries and $13$ minutes per precision.  Square shapes interpolate (log-log
between the measured sizes); rectangular shapes do not, and are either
calibrated for the shape ($118$ entries, a few seconds) or priced by
\emph{shape estimates} from the square calibration: GEMM-like entries at the
square size of equal flops, $\sqrt[3]{mkn}$ (Accelerate runs rectangular shapes
within 5\% of the square rate), and a level's passes at the square size of
equal traffic,
\begin{equation}\label{eq:equiv}
n_{\text{eq}} \;=\; \sqrt{\frac{(g_1g_2+r)\,mk + (g_2g_3+r)\,kn + (g_1g_3+r)\,mn}
{g_1g_2+g_2g_3+g_1g_3+3r}},
\end{equation}
the size at which a square problem moves the same words through the level's
passes \eqref{eq:traffic}.  The estimates predicted every measured rectangular
case within 2\% (Table~\ref{tab:rect}).

\paragraph{Sizes and shapes the levels do not divide.}
A plan runs at the smallest shape its levels divide: every dimension is
rounded up to the same multiple of $2^a3^b$ ($a\le10$, $b\le3$; 44 candidates),
the operands copied into zeroed matrices and the result copied back, priced at
$2\times$ the words at $80$~GB/s.  The planner takes the cheapest padded
candidate, so an $n$ that $3$ does not divide gets a $4\times4$ plan unless
padding to a multiple of $54$ pays.  Planning a new shape costs $0.07$~s in C
and $0.010$~s in Rust and is cached for $64$ shapes.

\paragraph{Validation.}
Predictions are within 0.3--3\% of the measured times at $n\ge8192$ in both
precisions throughout the development (first with the $4\times4$ menu, then
with streams, pipelines and the $3\times3\times6$ levels), and the planner's
choice was the fastest plan measured, or within 1\% of it, at every size and
precision tried, with two documented exceptions traced to noisy calibration
entries (a pipeline slope from two and four sub-problems; a $4096$ f32 run 15\%
slow throughout).  The classical variants themselves drift between runs
(13824 f64: $6.2$--$6.7$~s), so speedups carry about $\pm5$\%.

\section{Results}\label{sec:results}

\begin{table}[t]
\centering\small
\begin{tabular}{llrrrr}
\toprule
$n$, precision & plan & time (s) & over GEMM & relerr & err64\\
\midrule
6912 f32 & \lv{G} (Accelerate) & 0.195 & 1.00 & -- & $3.85\cdot10^{-6}$\\
 & \lv{D48/G} & 0.154 & 1.27 & $4.2\cdot10^{-6}$ & $2.81\cdot10^{-6}$\\
 & \lv{R336/G} & 0.164 & 1.19 & $3.9\cdot10^{-6}$ & $2.27\cdot10^{-6}$\\
 & \lv{R363/G} & 0.158 & 1.23 & $4.2\cdot10^{-6}$ & $1.76\cdot10^{-6}$\\
 & \lv{D48+4/D48/G6} & 0.149 & 1.31 & $4.3\cdot10^{-6}$ & $2.94\cdot10^{-6}$\\
 & \lv{R363+4/D48/G6} (auto) & \textbf{0.141} & \textbf{1.38} & $4.0\cdot10^{-6}$ & $2.26\cdot10^{-6}$\\
\midrule
13824 f32 & \lv{G} (Accelerate) & 1.565 & 1.00 & -- & $6.40\cdot10^{-6}$\\
 & \lv{D48/G} & 1.255 & 1.25 & $7.1\cdot10^{-6}$ & $4.16\cdot10^{-6}$\\
 & \lv{R336/G} & 1.244 & 1.26 & $6.7\cdot10^{-6}$ & $3.80\cdot10^{-6}$\\
 & \lv{R363/G} & 1.242 & 1.26 & $6.7\cdot10^{-6}$ & $2.73\cdot10^{-6}$\\
 & \lv{R336/R363/G} & 1.200 & 1.30 & $6.2\cdot10^{-6}$ & $3.34\cdot10^{-6}$\\
 & \lv{D48+4/D48/G6} & 0.949 & 1.65 & $7.0\cdot10^{-6}$ & $4.35\cdot10^{-6}$\\
 & \lv{R363+3/D48/G6} & 0.906 & 1.73 & $6.4\cdot10^{-6}$ & $3.49\cdot10^{-6}$\\
 & \lv{R363+4/R363/G6} (auto) & \textbf{0.907} & \textbf{1.73} & $6.6\cdot10^{-6}$ & $2.46\cdot10^{-6}$\\
\midrule
6912 f64 & \lv{G} (Accelerate) & 0.773 & 1.00 & -- & $6.69\cdot10^{-15}$\\
 & \lv{D48/G} & 0.612 & 1.26 & $7.8\cdot10^{-15}$ & $4.89\cdot10^{-15}$\\
 & \lv{R336/G} & 0.597 & 1.29 & $6.8\cdot10^{-15}$ & $4.53\cdot10^{-15}$\\
 & \lv{R363/G} & 0.647 & 1.19 & $6.8\cdot10^{-15}$ & $2.86\cdot10^{-15}$\\
 & \lv{R336/R363/G} & 0.594 & 1.30 & $8.3\cdot10^{-15}$ & $3.66\cdot10^{-15}$\\
 & \lv{D48+4/D48/G6} & 0.478 & 1.62 & $7.9\cdot10^{-15}$ & $5.60\cdot10^{-15}$\\
 & \lv{R363+3/R633/G6} (auto) & \textbf{0.451} & \textbf{1.71} & $6.7\cdot10^{-15}$ & $3.75\cdot10^{-15}$\\
\midrule
13824 f64 & \lv{G} (Accelerate) & 6.223 & 1.00 & -- & $1.23\cdot10^{-14}$\\
 & \lv{D48/G} & 4.970 & 1.25 & $1.40\cdot10^{-14}$ & $7.80\cdot10^{-15}$\\
 & \lv{R336/G} & 4.946 & 1.26 & $1.18\cdot10^{-14}$ & $7.21\cdot10^{-15}$\\
 & \lv{R363/G} & 4.919 & 1.27 & $1.30\cdot10^{-14}$ & $4.97\cdot10^{-15}$\\
 & \lv{R336/R363/G} & 3.703 & 1.68 & $1.45\cdot10^{-14}$ & $6.72\cdot10^{-15}$\\
 & \lv{D48+4/D48/G6} & 3.703 & 1.68 & $1.38\cdot10^{-14}$ & $8.28\cdot10^{-15}$\\
 & \lv{R363+3/D48/G6} & 3.467 & 1.80 & $1.25\cdot10^{-14}$ & $5.81\cdot10^{-15}$\\
 & \lv{SW/D48+2/D48/G6} & 3.407 & 1.83 & $1.85\cdot10^{-14}$ & $1.72\cdot10^{-14}$\\
 & \lv{SW/R363+3/R633/G6} (auto) & \textbf{3.251} & \textbf{1.91} & $1.51\cdot10^{-14}$ & $1.25\cdot10^{-14}$\\
\bottomrule
\end{tabular}
\caption{Square problems on the M4~Pro (\texttt{results/r336\_campaign2\_2026-10-28.txt};
the \lv{R363/G} rows and the $13824$ f64 \lv{R363+3/D48/G6} row from
\texttt{t42\_campaign} and \texttt{f64exact} of the same day): best of three
to five interleaved rounds; ``over GEMM'' is Accelerate's time divided by the
plan's; relerr is the max-norm difference from Accelerate's product; err64
the max-norm error against the exact product -- the f64 product in f32, the
compensated product in f64 (\texttt{results/f64exact\_2026-10-28.txt}) -- with
the GEMM's own error in the first row of each block.  ``auto'' is the
planner's choice.  Only err64 compares plans: the GEMM's own error is as
large as a level's, so relerr mixes the two (\lv{R336/G} and \lv{R363/G} at
$6912$ f64 have the same relerr and exact errors $4.53$ and
$2.86\cdot10^{-15}$).}
\label{tab:results}
\end{table}

Table~\ref{tab:results} compares, at two sizes and in both precisions, the
single levels, the best plan the $4\times4$ menu offered, and the planner's
mixed plan.  What the table says:

\begin{itemize}
\item As a single level the $3\times3\times6$ scheme ties the $4\times4$ scheme
  in time at $13824$ (in f32 at $6912$ it is 3--6\% slower, because the passes
  are a larger share there and its $L$ pass moves more) and is more accurate
  in every orientation: \lv{R336} by 7--19\%, \lv{R363} by 34--41\% in both
  precisions ($1.76$ against $2.81\cdot10^{-6}$ at $6912$ f32, $4.97$ against
  $7.80\cdot10^{-15}$ at $13824$ f64).  The planner knows the pass costs and
  uses the orientation whose cheapest pass faces the largest operand.
\item At matched time the mixed plans are the more accurate ones.  In each of
  the four cases the $3\times3\times6$-family plan that runs in the time of
  the best $4\times4$-only plan, \lv{D48+4/D48/G6}, has 23--43\% less error
  against the exact product and is 4--6\% faster: \lv{R363+4/D48/G6} at
  $6912$ f32 ($2.26$ against $2.94\cdot10^{-6}$), \lv{R363+4/R363/G6} at
  $13824$ f32 ($2.46$ against $4.35\cdot10^{-6}$), \lv{R363+3/R633/G6} at
  $6912$ f64 ($3.75$ against $5.60\cdot10^{-15}$) and \lv{R363+3/D48/G6} at
  $13824$ f64 ($5.81$ against $8.28\cdot10^{-15}$).  These plans, at
  $1.4$--$1.8\times$ over the GEMM, also have less error than a single
  $4\times4$ level at $1.25\times$ (\lv{D48/G}).  The first three are the
  planner's own choices, and its top six at $13824$ f32 are all
  $3\times3\times6$-family plans.
\item The exception is the planner's choice at $13824$ f64,
  \lv{SW/R363+3/R633/G6}: a Strassen--Winograd level on top buys $1.91\times$
  against $1.68$ for the best $4\times4$ plan at the price of 51\% more error
  ($1.25\cdot10^{-14}$ against $8.28\cdot10^{-15}$; the GEMM's own is
  $1.23\cdot10^{-14}$).  Against the $4\times4$ family's plan of the same
  design and time, \lv{SW/D48+2/D48/G6}, it is 5\% faster with 27\% less
  error, so the $3\times3\times6$ levels improve the trade at every point;
  but the planner optimizes time and takes the Strassen level whenever it
  pays, unless an error budget (\lv{auto:}$g$) forbids it.
\item Every plan without a Strassen level has less error than the GEMM's own
  in both precisions ($2.26$ against $3.85\cdot10^{-6}$ at $6912$ f32, $2.46$
  against $6.40\cdot10^{-6}$ at $13824$ f32, $3.75$ against
  $6.69\cdot10^{-15}$ at $6912$ f64, $5.81$ against $1.23\cdot10^{-14}$ at
  $13824$ f64): a fast multiplication that is also the more accurate one
  (Section~\ref{sec:accuracy} says why).  The difference from Accelerate's
  product, the only f64 error the earlier campaigns recorded, hides all of
  this: it mixes the GEMM's error with the plan's.
\item The three-orientation cycle alone (\lv{R336/R363/R633/G}) does not pay at
  these sizes: its leaves of $n/54$ are too small for the GEMM.  Depth pays
  through the pipelined two-level plans.
\end{itemize}

The planner's choices by size, from the calibration at $13824$: f32 \lv{G} at
$1728$, \lv{R363/G6} at $3456$, \lv{R363+4/D48/G6} at $6912$,
\lv{R363+4/R363/G6} at $13824$; f64 \lv{R363/G6} at $1728$, \lv{SW+3/R363/G6}
at $3456$, \lv{R363+3/R633/G6} at $6912$, \lv{SW/R363+3/R633/G6} at $13824$.
Sizes that $3$ does not divide keep their $4\times4$ plans (\lv{D48+4/D48/G6}
at $8192$ f32, $1.40\times$; \lv{D48+3/D48/G6} at $9000$ f64 padded to $9008$,
$1.69\times$).

\begin{table}[t]
\centering\small
\begin{tabular}{llllrr}
\toprule
$m\times k\times n$ & calibration & prec. & plan & predicted / measured (s) & over GEMM\\
\midrule
$3456\times3456\times1728$ & for the shape & f32 & \lv{R363/G6} & 0.0117 / 0.0115 & 1.05\\
 & & f64 & \lv{R363/G6} & 0.0366 / 0.0367 & 1.30\\
$1000\times3000\times2000$ & estimated & f32 & \lv{G} & -- & 1.00\\
 & & f64 & \lv{SW/G4} & 0.0133 / 0.0144 & 1.04\\
$2048\times4096\times1024$ & estimated & f32 & \lv{G} & -- & 1.00\\
 & & f64 & \lv{D48/G6} & 0.0177 / 0.0180 & 1.10\\
$6912\times6912\times3456$ & estimated & f32 & \lv{SW+4/D48/G6} & 0.0790 / 0.0796 & 1.26\\
 & & f64 & \lv{D48+4/R363/G6} & 0.2578 / 0.2565 & 1.52\\
\bottomrule
\end{tabular}
\caption{Rectangular problems through the library's entry points
(\texttt{results/\allowbreak rect\_\allowbreak api\_\allowbreak 2026-10-28.txt}): the planner's plan, its predicted
and measured time, and the speedup over Accelerate's GEMM on the same shape.
``Estimated'' shapes are priced from the square calibration by
\eqref{eq:equiv}.}
\label{tab:rect}
\end{table}

Rectangular problems (Table~\ref{tab:rect}) go through the same planner over
$(m,k,n)$: a level applies when its grid divides the current shape, so a
$6912\times6912\times3456$ product gets a $4\times4$ level whose products are
$1728\times1728\times864$ and then a $3\times6\times3$ level in f64.  The shape
estimates are within 2\% of the measurements wherever a level pays, and where
none does the planner correctly returns the GEMM (the first version of the
estimates had none, and padded $1000\times3000\times2000$ to $3072^3$, 30
times slower than Accelerate).

\subsection{The GPU as the alternative}\label{sec:gpu}

The other engine on the machine is the 20-core GPU.  We measured Apple's own
GEMM paths on it (\path{gpu_gemm.m}, Section~\ref{sec:glossary} for the
names): MPSGraph's matrix multiplication, the faster of the two, and the
older \lv{MPSMatrixMultiplication}, within 3\% of it at $13824$ and 10\%
slower at $3456$; in f32, since Metal, MPS and MPSGraph have no 64-bit
floating type.  The inputs are the benchmark's, wrapped as shared Metal
buffers without a copy (the memory is unified), the time is the wall clock
from encoding to completion, five rounds interleaved with Accelerate's GEMM on
the same inputs, the error against the f64 product.  Neither MLX nor PyTorch
was installed, so these are the vendor's GEMMs, as Accelerate is on the CPU
side.

\begin{table}[t]
\centering\small
\begin{tabular}{lrrrrlrr}
\toprule
 & \multicolumn{2}{c}{GPU, MPSGraph} & Accelerate & GPU and Acc. & \multicolumn{3}{c}{best mixed plan}\\
\cmidrule(lr){2-3}\cmidrule(lr){6-8}
$n$ (f32) & time (s) & TFLOP/s & time (s) & error & plan & time (s) & error\\
\midrule
3456 & 0.013 & 6.3 & 0.025 & $3.04\cdot10^{-6}$ & \lv{R363/G6} & 0.021 & $1.28\cdot10^{-6}$\\
6912 & 0.099 & 6.6 & 0.195 & $3.85\cdot10^{-6}$ & \lv{R363+4/D48/G6} & 0.141 & $2.26\cdot10^{-6}$\\
13824 & 0.946 & 5.6 & 1.565 & $6.40\cdot10^{-6}$ & \lv{R363+4/R363/G6} & 0.907 & $2.46\cdot10^{-6}$\\
\bottomrule
\end{tabular}
\caption{The GPU as a GEMM (\texttt{results/gpu\_gemm\_2026-10-28.txt}): best
of five rounds; the GPU's f32 error equals Accelerate's, hence one column.
With f16 inputs and an f32 result the GPU takes $0.011$, $0.090$ and
$0.790$~s with errors $2.5$--$2.8\cdot10^{-4}$.  Accelerate's and the plans'
numbers are the campaigns' (Table~\ref{tab:results}, and the $3456$ exact
errors of \texttt{results/f64exact\_2026-10-28.txt}).}
\label{tab:gpu}
\end{table}

What Table~\ref{tab:gpu} says:
\begin{itemize}
\item The GPU's f32 product is bit for bit Accelerate's at every size.  Both
  accumulate each output in increasing $k$ with fused multiply-adds and
  neither splits the inner dimension, so both equal the sequential FMA chain
  (64 sampled outputs recomputed as that chain agree with both), and the
  GPU's error is the GEMM's.  The mixed plans have $1.7$--$2.6\times$ less
  error than either.
\item In time the GPU is $1.6\times$ ahead of the best plan at $3456$ and
  $1.4\times$ at $6912$, and tied at $13824$, where its throughput falls to
  $5.6$~TFLOP/s from $6.6$ at $6912$ and the CPU plan's $0.907$~s edges its
  $0.946$.
\item The GPU's f16-input path (f32 result) is only 14\% faster than its f32
  at $13824$, and its error, $2.5$--$2.8\cdot10^{-4}$, is the rounding of the
  inputs to 11 bits.
\item In f64 the GPU does not compete, and the plans' $1.71$--$1.91\times$
  over Accelerate's dgemm stands alone.
\item The two engines share the package's power and thermal budget:
  Accelerate's sgemm interleaved with GPU runs took $1.81$--$1.88$~s at
  $13824$ against $1.565$ alone.  A hybrid that keeps both busy has to
  reckon with that; by estimate a $3\times3\times6$ level over GPU leaves at
  $13824$ f32 would run in about $0.77$~s ($40/54$ of the GPU's $0.946$ plus
  $0.07$~s of CPU passes), $1.2\times$ over the GPU alone, if the leaf
  products keep the GPU's throughput.  Not measured.
\end{itemize}

\section{The ternary trial: cheap passes against accurate ones}\label{sec:ternary}

Perminov's catalogue \cite{Perminov,Perminov25} lists, beside the ranks, the
schemes with coefficients in $\{-1,0,1\}$ whose additive complexity his
parallel common-subexpression search reduced; for $3\times3\times6$ it has four
rank-$42$ ternary schemes reduced to $139$, $140$, $145$ and $150$ additions
for the three linear forms together (file
\texttt{3x3x6\_m42\_cr139\_cn287\_ZT\_reduced.json} and its siblings; the
naive counts are $287$--$369$).  Against DPS's accurate rank-$40$ scheme,
whose nine optimized programs take $374$, $383$ and $353$ operations per
orientation, this is the opposite design: two products more (saving $0.778$
against $0.741$), two more streams per pass ($51/60/60$), and a third of the
pass arithmetic.  Is the arithmetic worth anything once the passes are fused,
and what does the accuracy cost?

\paragraph{The scheme.}
The reduced JSON is a straight-line program already (fresh variables for the
common subexpressions, then the forms); each of the four schemes was expanded
to its matrices and checked exactly against the Brent equations.  The
$139$-addition scheme is the most accurate of the four by far (amplification
ratio $5.07$ in the $3\times3\times6$ orientation against $7.6$--$9.5$ for the
others), so it is the one tried; Table~\ref{tab:levels} gives its three
orientations.  Its programs are Perminov's forms for the three maps the
$3\times3\times6$ orientation needs ($L$: $33$ additions, $R$: $39$, $P$: $67$),
and for the other orientations the same forms re-indexed where an
orientation reuses a map and transposed by PLinOpt's \texttt{transpozer}
where it needs the transposed map ($3\times6\times3$: $39$, $43{+}6$ negations,
$66$; $6\times3\times3$: $43{+}6$, $33$, $63$), all nine checked by
\texttt{SLPchecker} against the matrices.  PLinOpt's optimizer (the fork's
\texttt{-s 1 -Q -O 500}) on the same nine matrices finds $34$, $41$, $71$,
$41$, $45$, $67$, $45$, $34$, $66$ additions: one to four more than Perminov's
reduction on every matrix, in $0.2$~s each.  The levels \lv{T336}, \lv{T363},
\lv{T633} went into the C library beside the R levels, outside the planner's
menu, and every plan containing them computes the product (3456, both
precisions, including pipelined plans).

\begin{table}[t]
\centering\footnotesize
\begin{tabular}{llrrlrrl}
\toprule
$n$, prec. & plan & R (s) & T (s) & T/R & R error & T error & T/R\\
\midrule
13824 f32 & \lv{X336/G} & 1.228 & 1.286 & 1.05 & 3.80e-6 & 1.28e-5 & 3.4\\
 & \lv{X363/G} & 1.242 & 1.304 & 1.05 & 2.73e-6 & 6.25e-6 & 2.3\\
 & \lv{X633/G} & 1.276 & 1.339 & 1.05 & 3.95e-6 & 1.01e-5 & 2.6\\
 & \lv{X336/X363/X633/G} & 2.538 & 2.874 & 1.13 & 4.95e-6 & 3.53e-5 & 7.1\\
 & \lv{X363+4/D48/G6} & 0.909 & 0.954 & 1.05 & 3.49e-6 & 6.23e-6 & 1.8\\
 & \lv{X363+4/X363/G6} & 0.906 & 0.994 & 1.10 & 2.46e-6 & 9.04e-6 & 3.7\\
\midrule
13824 f64 & \lv{X336/G} & 4.858 & 5.089 & 1.05 & 1.18e-14 & 2.62e-14 & 2.2\\
 & \lv{X363/G} & 4.919 & 5.132 & 1.04 & 1.30e-14 & 1.34e-14 & 1.0\\
 & \lv{X336/X363/X633/G} & 7.546 & 8.763 & 1.16 & 1.24e-14 & 8.08e-14 & 6.5\\
 & \lv{SW/X363+3/X633/G6} & 3.236 & 3.573 & 1.10 & 1.51e-14 & 5.61e-14 & 3.7\\
 & \lv{X363+3/X633/G6} & 5.624 & 6.188 & 1.10 & 1.26e-14 & 2.63e-14 & 2.1\\
\midrule
6912 f32 & \lv{X336/G} & 0.165 & 0.173 & 1.05 & 2.27e-6 & 8.00e-6 & 3.5\\
 & \lv{X363/G} & 0.158 & 0.160 & 1.01 & 1.76e-6 & 3.87e-6 & 2.2\\
 & \lv{X363+4/D48/G6} & 0.142 & 0.147 & 1.04 & 2.26e-6 & 4.07e-6 & 1.8\\
 & \lv{X363+3/X633/G6} & 0.160 & 0.173 & 1.08 & 2.20e-6 & 8.72e-6 & 4.0\\
\midrule
6912 f64 & \lv{X336/G} & 0.598 & 0.626 & 1.05 & 6.78e-15 & 1.47e-14 & 2.2\\
 & \lv{X363+4/D48/G6} & 0.453 & 0.478 & 1.06 & 7.03e-15 & 8.81e-15 & 1.3\\
 & \lv{X363+3/X633/G6} & 0.451 & 0.497 & 1.10 & 6.74e-15 & 1.63e-14 & 2.4\\
\bottomrule
\end{tabular}
\caption{The ternary $\mm{3\times3\times6:42}$ levels (T) against the accurate
$\mm{3\times3\times6:40}$ levels (R) on the same plans
(\texttt{results/t42\_campaign\_2026-10-28.txt}; \lv{X} stands for \lv{R} or
\lv{T}): best of five (f32) or three (f64) interleaved rounds, seconds; the
f32 error is against the f64 product, the f64 error against Accelerate's
product, which in f64 includes Accelerate's own error, as large as an R
level's (Section~\ref{sec:results}), so the f32 ratios carry the
comparison.}
\label{tab:ternary}
\end{table}

\paragraph{Measured.}
Table~\ref{tab:ternary} answers both questions.  The pass arithmetic is worth
nothing: the T levels' passes take exactly the R levels' times (at $13824$
f32, $L/R/P$ $0.036/0.020/0.017$~s against $0.034/0.019/0.017$; at $6912$,
$0.009/0.005/0.005$ for both), because a fused pass is bound by the blocks it
moves and the T level moves two more per pass.  The two extra products cost
what the rank says: the GEMM part of a single level is 5--7\% longer, a
single T level is 4--5\% slower than its R counterpart at every size and
precision (the one tie, \lv{T363/G} at $6912$ f32, is within the run's
noise), a two-level T plan 8--10\% slower, and the three-orientation cycle
13--16\% slower.  The accuracy cost is the amplification: $2.2$--$3.5\times$ the
R level's error as a single level, $3.7$--$4\times$ in two-level plans and
$6.5$--$7\times$ in the cycle, in line with the ratio of the amplification
factors ($5.07/2.36 = 2.1$ for the $3\times3\times6$ orientation, the
measured $2.2$--$3.5$ including the passes' own roundings).  In f32 every T
plan's error is above the GEMM's own (\lv{T336/G} $1.28\cdot10^{-5}$ against
$6.40\cdot10^{-6}$ at $13824$) where every R plan's is below it.  The
planner's menu was left as it was: a level that is slower and less accurate
than one already there on every plan cannot be chosen.

\paragraph{What it shows.}
The trial is the cleanest confirmation of the cost model's premise.  Two
schemes of the same shape, with the same pass structure, differ in rank by
two and in pass arithmetic by a factor of $2.7$; the measured difference is
the rank's, to the percent, and none of the arithmetic's.  For a fused-pass
library the objectives of a scheme search are rank and amplification, in
that order, and additive complexity is not among them.  The caveat is the
one stated in Section~\ref{sec:levels}: arithmetic shows when a program is
wasteful enough to make a pass compute-bound, as the $18$-operation-per-output
direct sums were; Perminov's $139$ additions and DPS's optimized programs are
both well inside the bandwidth-bound regime.  Whether the class of the
ternary scheme contains an accurate point is a separate question, and the
answer is no: over $GL_3\times GL_3\times GL_6$ the class's amplification
ratio is at least
$2.825$ in the $3\times3\times6$ orientation and $2.831$ in $3\times6\times3$
(eight starts each), the $4\times4$ class's $2\sqrt2$ and 26\% above the
accurate $\mm{3\times3\times6:40}$ class's $2.242$; the $139$-addition point
sits $1.8\times$ above its own floor, and its $\ginf$ floors ($5.12$ and $4.82$)
are likewise above the accurate class's ($6.37$ and $3.90$) in the orientation
that matters.  Even at its best point the class would match the $4\times4$
scheme's accuracy with two products more than the accurate $3\times3\times6$
scheme.

\section{Two implementations}\label{sec:impl}

The library is C (\texttt{fmm.h}, \texttt{libfmm.a}: \texttt{fmm\_sgemm},
\texttt{fmm\_dgemm} and their \texttt{\_mkn} forms for an $m\times k$ times
$k\times n$ product with separate leading dimensions, \texttt{fmm\_calibrate},
\texttt{fmm\_plan}, \texttt{fmm\_release}) over Accelerate, with the kernels
generated from the schemes' programs by \texttt{gen\_kernels.py}: every
program in PLinOpt's syntax becomes one fused loop in f32 and in f64, plus an
f64-arithmetic variant for f32 blocks.  Plans are made once per (shape,
precision, budget) and kept; the workspace is kept between calls; a call
outside the calibrated range or without a calibration runs the classical GEMM,
so a call is always correct.  The benchmark driver includes the library, so
both run the same code, and every plan prints its predicted time next to the
measured one.

The Rust port (crate \texttt{fmm}, \texttt{rust/}) is the same library module
for module -- plans and shapes, the calibration in the C file format so that
either library reads the other's cache, the planner with pipelines, budget
and padding, the executor on a rayon pool, the entry points -- with the
fused passes as safe Rust emitted by the same generator (each slice a
separate parameter, cut to the row length first: the compiler then
vectorizes as \texttt{restrict} and the pragma let it in C).  The only unsafe
code is the three FFI declarations for Accelerate behind safe wrappers that
check the slices.  Cross-checked against the C library on the shared
calibration, 26 cases (square sizes from $1000$ to $9000$, a leading dimension
above $n$, an error budget, four rectangular shapes, both precisions) give the
same plans, the same predicted times to the digit and results bit for bit.
Speeds match at large sizes (f64 $8192$: $0.801$ against $0.815$~s; f32
$6912\times6912\times3456$: $0.0789$ against $0.0796$); the Rust passes on
blocks under about $600$ are $1.2$--$2\times$ slower (dispatch and per-chunk
allocation), which only shows when Rust calibrates for itself.  Two traps
found on the way: a zeroed \texttt{Vec} gets its pages on first write, so
timed buffers must be written first; and a lone multi-threaded GEMM issued
right after GEMM streams on the rayon pool runs 15--25\% slower and stays so,
which the leaves' timing avoids by running where the library runs a
classical plan.

\section{Discussion}\label{sec:disc}

The mixed approach rests on two measurements that contradict the usual
accounting.  A fused level's cost is the blocks it moves, so the right
measure of a scheme is its saving per block moved and its stream count, not
its rank or the length of its programs; and the accuracy of a level is set by
the point in its orbit and, at depth, by the roundings of its $P$ program,
not by the operation count.  Under these measures the menu is short: two
accurate schemes and Strassen--Winograd for the top of deep f64 plans.  A
scheme enters the menu only if a plan containing it beats every plan without
it on this machine, and the planner decides that per shape and precision from
measurements, which is why the library can take a new level (or a new
machine) without any change but a calibration.

The limits are also measured.  The GEMM is 85--95\% of every plan, so no menu
can take more than the products it saves, and below $n\approx3000$ nothing
pays in f32 ($1728$ in f64).  A level's passes must not fall behind the
streams it feeds: six streams of $1024$-blocks in f64 are the most sensitive
to pass traffic, and pipelining gains nothing there.  Shapes above
$4\times4$ cannot help because of the stream count, whatever their rank, and
Kronecker composites because they inherit their pieces' amplification and
multiply their streams.  And the ternary trial (Section~\ref{sec:ternary})
closes the remaining door on the cheap side: when the passes are already
bandwidth-bound, a program with a third of the additions buys nothing, while
two extra products and $2.8\times$ the amplification cost speed and accuracy
at every size.

What is next.  A dyadic point of the $3\times6\times3$ orientation closer to
its class's amplification floor (DPS's is 14\% above it) would improve the
plans that use it; the tools exist from the $4\times4$ search.  Fusing the $L$
and $R$ passes into GEMM operand packing, the ``Strassen reloaded'' design,
needs an open GEMM kernel; on this machine the vendor's SME kernels are
closed, so it waits for another machine (x86-64 with BLIS, or the Apple GPU).
The planner's $0.07$~s per new shape in C should be a millisecond.  The
error model should be fit to the exact product instead of Accelerate's, whose
own error it now includes (Section~\ref{sec:accuracy}); with that done the
planner's default objective, time, could become time at the GEMM's own
accuracy, which every plan without a Strassen level meets.  The
library takes $C=AB$ only: $\alpha$, $\beta$ and transposes are a day's work.
And the Rust passes on small blocks should be brought to the C library's
speed.

\section{Glossary: the machine, its software, and the report's terms}\label{sec:glossary}

The report uses Apple's names for the hardware and the libraries, and a few
of Arm's.  In one place, followed by the report's own terms and abbreviations:

\begin{description}
\item[Apple silicon, M4~Pro.]  Apple's system on a chip for Macs.  The M4~Pro
  of this report has 14 CPU cores in three clusters (two of five performance
  cores, one of four efficiency cores), a 20-core GPU, and 64~GB of memory
  that all of them share.  A cluster is Apple's group of
  cores around one level-2 cache; each performance cluster of the M4 also
  has one matrix unit.
\item[Performance and efficiency cores.]  The two kinds of CPU core.  The
  performance cores run the passes and the GEMM threads here; the efficiency
  cores are slower and share a smaller cache, and the thread counts of
  Section~\ref{sec:machine} are chosen so that timed work stays off them.
\item[SME, Scalable Matrix Extension.]  Arm's instruction set extension for
  matrix arithmetic, implemented in the M4 as a matrix coprocessor with a
  large tile register and an outer-product instruction that multiplies a
  column of $A$ by a row of $B$ and accumulates into the tile.  The M4~Pro has
  one SME unit per performance cluster, shared by that cluster's cores; it
  succeeds Apple's earlier undocumented AMX unit.  Accelerate's GEMM runs on
  it: a GEMM call keeps about two cores busy and reaches $3.3$~TFLOP/s in
  f32, far above what the cores' vector units could give, with the GPU's
  utilization counter at zero.  It accumulates in the order of the
  sequential FMA chain, which is why its result equals the GPU's bit for
  bit.  Its kernels are Apple's and closed, which keeps the ``Strassen
  reloaded'' design of fusing the passes into the GEMM's operand packing off
  this machine.
\item[NEON.]  Arm's 128-bit SIMD vector instruction set, on every core.
  GEMMs on the idle cores with it were tried and not adopted
  (Section~\ref{sec:machine}).
\item[Accelerate.]  Apple's numerical library: BLAS and LAPACK, vDSP and
  more, with the kernels for Apple silicon.  Its BLAS is the vendor GEMM of
  this report, called through the C interface \lv{cblas\_sgemm} and
  \lv{cblas\_dgemm}; the threading inside a call is Accelerate's own, and
  the library's streams are several single-threaded calls side by side.
\item[BLAS, GEMM, sgemm, dgemm.]  The Basic Linear Algebra Subprograms
  interface; GEMM is its general matrix multiply $C \leftarrow \alpha AB +
  \beta C$, sgemm in single precision (f32), dgemm in double (f64).  ``The
  GEMM'' in this report is Accelerate's.
\item[FMA, fused multiply-add.]  The instruction computing $a\cdot b + c$
  with a single rounding; every modern GEMM is built from it.  One output
  accumulated as $c \leftarrow \mathrm{fma}(a_k, b_k, c)$ for $k$ in order is
  the sequential FMA chain that both Accelerate's SME kernels and Apple's GPU
  kernels compute.
\item[f32, f64, f16, bf16.]  IEEE single and double precision (24- and
  53-bit significands), IEEE half precision (11 bits) and bfloat16 (8 bits).
  The CPU paths of this report are f32 and f64; Apple's GPU has f32, f16 and
  bf16 and no f64.
\item[GPU, unified memory.]  The M4~Pro's 20-core graphics processor,
  programmed through Metal.  It reads the same memory as the CPU, so a CPU
  array is handed to it without a copy by wrapping it as a shared Metal
  buffer; the GPU numbers of Section~\ref{sec:gpu} include no transfers.
\item[Metal.]  Apple's GPU programming interface and shading language, the
  counterpart of CUDA or Vulkan: devices, command queues, command buffers,
  compute kernels, buffers.
\item[MPS, Metal Performance Shaders.]  Apple's library of ready GPU kernels
  on top of Metal, among them \lv{MPSMatrixMultiplication}, a GEMM on Metal
  buffers.
\item[MPSGraph.]  Metal Performance Shaders Graph: Apple's framework for
  tensor computation on the GPU, in which a graph of operations (here two
  placeholders and one matrix multiplication) is compiled to Metal kernels
  and run.  It is what Core ML and PyTorch's Apple backend use underneath;
  its GEMM is the faster of Apple's two.
\item[Grand Central Dispatch.]  Apple's thread-pool interface (libdispatch);
  \lv{dispatch\_apply} runs a block over an index range on the pool.
  The library uses both it and POSIX threads; the benchmark's compensated reference product runs on the pool through \lv{dispatch\_apply}.
\item[Apple clang, macOS~26.]  The compiler and operating system of the
  measurements.  The C library is C11 with Accelerate; the GPU tool is
  Objective-C, because the GPU frameworks have no C interface.
\item[TFLOP/s.]  $10^{12}$ floating-point operations per second, counting
  $2mkn$ for a product whatever the algorithm, so a plan's rate is the
  classical count divided by its time.
\end{description}

\begin{description}
\item[GEMM.]  See above: the general matrix multiply of BLAS, here
  Accelerate's, the leaf of every plan and the baseline of every speedup.
\item[Scheme $\mm{g_1\times g_2\times g_3:r}$, rank.]  A bilinear algorithm
  that multiplies a $g_1\times g_2$ by a $g_2\times g_3$ block matrix with $r$
  block products instead of $g_1g_2g_3$; $r$ is its rank.  Its three matrices
  $L$, $R$, $P$ are Section~\ref{sec:levels}'s.
\item[Level, pass, plan, leaf.]  A level is one application of a scheme to
  the current problem, executed as three fused passes ($L$, $R$, $P$) around
  its products; a plan is the sequence of levels from the top down to the
  leaf, the GEMM; the notation is in Section~\ref{sec:intro}.
\item[\lv{D48}, \lv{R336}, \lv{R363}, \lv{R633}, \lv{SW}, \lv{T336},
  \lv{T363}, \lv{T633}, \lv{X}.]  The levels of the menu (Table~\ref{tab:levels}):
  the accurate $4\times4$ scheme, the accurate $3\times3\times6$ scheme in
  its three orientations, Strassen--Winograd, and the ternary
  $3\times3\times6$ scheme of Section~\ref{sec:ternary}; \lv{X} in
  Table~\ref{tab:ternary} stands for \lv{R} or \lv{T}.
\item[\lv{G}, \lv{G}$k$, \lv{B2}, \lv{+p}.]  The leaves and the pipeline
  mark: \lv{G} is one Accelerate GEMM with its own threading, \lv{G}$k$ the
  products as $k$ single-threaded GEMM streams, \lv{B2} the classical
  $2\times2$ blocked product as a level (eight products, the calibration's
  reference); \lv{+p} after a level runs its passes on $p$ threads overlapped
  with the products below.
\item[Streams, pipelines.]  Section~\ref{sec:machine}: several
  single-threaded GEMM calls side by side, and a level's passes overlapped
  with the products of the level below.
\item[err64, relerr.]  The max-norm relative error against the exact
  product (Section~\ref{sec:intro}), and the same quantity against
  Accelerate's product, which includes the GEMM's own error and is kept only
  for comparison with the earlier campaigns.
\item[Error growth, \lv{max\_growth}, \lv{auto:}$g$.]  A plan's predicted
  error divided by the GEMM's, the product of the model's per-level factors
  $f_X$; the budget $g$ a call sets on it, with $0$ meaning no budget
  (Section~\ref{sec:planner}).
\item[$\amp$, $\ginf$, $N_E$, $e_0$.]  The amplification measures of a
  scheme (root-sum-square and worst-case, Section~\ref{sec:accuracy}); the
  size at which the error model is calibrated ($2304$) and its classical
  error scale there.
\item[Kronecker composite.]  The scheme obtained as the Kronecker product of
  two schemes (one level of each, fused), screened out in
  Section~\ref{sec:menu}.
\item[DPS.]  Dumas, Pernet and Sedoglavic, the authors of the accurate
  schemes and of PLinOpt \cite{DPS24,DPS25,DPS26a,DPS26b}; DPSST adds
  Stapleton and Tichavsk\'y \cite{DPSST26}.
\item[PLinOpt, SLP, SLPchecker.]  DPS's library of tools for linear
  programs \cite{PLinOpt}: a straight-line program (SLP) is a list of
  assignments computing a linear map; PLinOpt's optimizer shortens one, its
  transpozer transposes one, and SLPchecker verifies one against its matrix,
  which every program in this report passed.
\item[Orbit, $GL_n$.]  The schemes equivalent to a given one under the
  general linear group $GL_n$ of invertible $n\times n$ matrices acting on
  the three sides; the orbit floors of Section~\ref{sec:accuracy} are minima
  over it, found by L-BFGS, a quasi-Newton minimizer.
\item[RMS, IEEE 754, SIMD.]  Root mean square; the floating-point
  standard (f32, f64, f16); single instruction, multiple data (vector
  instructions).
\item[POSIX threads, FFI, JSON, C11.]  The standard thread interface the C
  library uses; the foreign-function interface through which Rust calls
  Accelerate; the text format of Perminov's catalogue; the 2011 C standard.
\item[BLIS, CUDA, Vulkan, Core ML, PyTorch, MLX.]  Named for comparison only:
  an open BLAS framework (for the x86-64 machine the design of
  Section~\ref{sec:disc} waits for), Nvidia's and Khronos's GPU interfaces,
  Apple's and Meta's machine-learning frameworks, and Apple's array
  framework.
\item[GB/s, TFLOP/s.]  Gigabytes per second of memory traffic; $10^{12}$
  floating-point operations per second as defined above.
\end{description}

\section{Data and code}\label{sec:data}

Everything is in the repository \cite{logic}, directory
\path{matmul/fastmm/}: the library (\path{fmm.c}, \path{fmm_impl.h},
\path{fmm.h}), the kernel generator (\path{gen_kernels.py}), the programs
of the $3\times3\times6$ levels (\path{schemes/opt/r*}, PLinOpt-optimized;
\path{schemes/r*}, the direct sums) and of the ternary levels
(\path{schemes/t*}), the Rust crate (\path{rust/}), the benchmark and the
test (\path{bench_fmm.c}, \path{test_fmm.c}), the calibrations
(\path{results/calib5_f32.txt}, \path{calib5_f64.txt}) and the raw
results of every run quoted here:
\path{results/r336_campaign_2026-10-28.txt},
\path{r336_campaign2_2026-10-28.txt} (Table~\ref{tab:results}),
\path{r336_opt_vs_direct_2026-10-28.txt},
\path{rect_api_2026-10-28.txt} (Table~\ref{tab:rect}),
\path{t42_campaign_2026-10-28.txt} (Table~\ref{tab:ternary}),
\path{f64exact_2026-10-28.txt} (the f64 err64 column of
Table~\ref{tab:results}), \path{gpu_gemm_2026-10-28.txt} (Table~\ref{tab:gpu};
the tool is \path{gpu_gemm.m}, built by \texttt{make build/gpu\_gemm}, the
script \path{run_gpu.sh}), and the
earlier campaigns of October 24--26 (streams, pipelines, accuracy, the
$4\times4$ programs).  The menu screen is \path{matmul/levels/}
(\path{levels.py}, \path{levels.csv}, \path{orbit_rect.py},
\path{cost_rect.py}, README).  The accuracy study of the $4\times4$ class
is \path{doc/matmul_4x4_rank48_accurate262.pdf} \cite{S26b}.  The running
log of the track, with every measurement in the order it was made, is
\path{doc/fastmm_plan.md}.

{\small
\begin{thebibliography}{99}
\bibitem{DPS24} J.-G.~Dumas, C.~Pernet, A.~Sedoglavic, \emph{Strassen's
algorithm is not optimally accurate}, Proc.\ ISSAC 2024, 254--263.
\bibitem{DPS25} J.-G.~Dumas, C.~Pernet, A.~Sedoglavic, \emph{A
non-commutative algorithm for multiplying $4\times4$ matrices using
$48$ non-complex multiplications}, arXiv:2506.13242 (2025).
\bibitem{DPS26a} J.-G.~Dumas, C.~Pernet, A.~Sedoglavic, \emph{A more
accurate rational non-commutative algorithm for multiplying $4\times4$
matrices using $48$ multiplications}, arXiv:2603.18699 (v1 19~March
2026, v3 29~July 2026).
\bibitem{DPS26b} J.-G.~Dumas, C.~Pernet, A.~Sedoglavic, \emph{Towards
automated generation of fast and accurate algorithms for recursive
matrix multiplication}, J.\ Symbolic Comput.\ 134 (2026) 102524.
\bibitem{DPSST26} J.-G.~Dumas, C.~Pernet, A.~Sedoglavic, J.~Stapleton,
P.~Tichavsk\'y, \emph{Fast matrix multiplication via recursive
$\langle4\times4\times4:48\rangle$ algorithms into practice},
arXiv:2609.12027 (v1 10~September 2026, v2 2~October 2026).
\bibitem{PLinOpt} J.-G.~Dumas, B.~Grenet, C.~Pernet, A.~Sedoglavic,
\emph{PLinOpt, a collection of C++ routines handling linear \&
bilinear programs}, \url{https://github.com/jgdumas/plinopt}; data
files \texttt{4x4x4\_48\_accurate\_262} (commit \texttt{b441528},
13~October 2026), \texttt{3x3x6\_40\_DPS-accurate}, \texttt{3x6x3\_40\_DPS-accurate},
\texttt{6x3x3\_40\_DPS-accurate}; tools \texttt{optimizer}, \texttt{transpozer},
\texttt{SLPchecker}.
\bibitem{HSvdG16} J.~Huang, T.~M.~Smith, G.~M.~Henry, R.~A.~van de Geijn,
\emph{Strassen's algorithm reloaded}, Proc.\ SC~'16 (International
Conference for High Performance Computing, Networking, Storage and
Analysis), 2016.
\bibitem{KW} M.~Kauers, J.~Wood, \emph{Consequences of the Moosbauer--Poole
algorithms}, preprint (2025).
\bibitem{MP25} J.~Moosbauer, M.~Poole, \emph{Flip graphs with symmetry and
new matrix multiplication schemes}, preprint (2025).
\bibitem{MSW25} E.~M\r{a}rtensson, P.~Stankovski Wagner, \emph{The number of
the beast: reducing additions in fast matrix multiplication algorithms
for dimensions up to 666}, Proc.\ ACDA 2025; the library
\texttt{fmm\_add\_reduction}, \url{https://github.com/werekorren/fmm_add_reduction}
(commit \texttt{951cd93}, 8~May 2025), whose \texttt{algorithms/} directory
distributes the Kauers--Wood, Moosbauer--Poole--Kemper and AlphaEvolve
schemes used in the screen.
\bibitem{Novikov25} A.~Novikov et al., \emph{AlphaEvolve: a coding agent
for scientific and algorithmic discovery}, arXiv:2506.13131 (2025).
\bibitem{Perminov} A.~I.~Perminov, \emph{FastMatrixMultiplication}: a
catalogue of fast matrix multiplication schemes,
\url{https://github.com/dronperminov/FastMatrixMultiplication}, commit
\texttt{3ee3284} (28~October 2026); file
\texttt{schemes/results/addition\_reduced\_ZT/3x3x6\_m42\_cr139\_cn287\_ZT\_reduced.json}.
\bibitem{Perminov25} A.~I.~Perminov, \emph{Parallel heuristic exploration
for additive complexity reduction in fast matrix multiplication},
arXiv:2512.13365 (2025).
\bibitem{Smirnov13} A.~V.~Smirnov, \emph{The bilinear complexity and
practical algorithms for matrix multiplication}, Comput.\ Math.\ Math.\
Phys.\ 53 (2013) 1781--1795.
\bibitem{S26} G.~Sidebottom, \emph{201 operations for rank-48 $4\times4$
matrix multiplication over rings containing $1/2$}, preprint (2026);
PLinOpt data file \texttt{4x4x4\_48\_201}.
\bibitem{S26b} G.~Sidebottom, \emph{The 262-operation
$\langle4\times4\times4:48\rangle$ algorithm is also the more accurate one:
error-bound exponent $9/4$ and about a quarter less error}, draft of
27~October 2026, \texttt{doc/matmul\_4x4\_rank48\_accurate262.pdf} in \cite{logic}.
\bibitem{logic} G.~Sidebottom, \emph{logic}, source repository,
\url{https://github.com/gsidebottom/logic} (2026).
\end{thebibliography}
}

\end{document}
